\documentclass[11pt]{article}
\usepackage[margin=1.15in]{geometry}
\usepackage{amsmath,amssymb}
\usepackage[hidelinks]{hyperref}
\pagestyle{empty}
\setlength{\parskip}{5pt}
\begin{document}

\noindent
Claes Johnson\\
KTH Royal Institute of Technology\\
SE-100 44 Stockholm, Sweden\\
\texttt{claesjohnson@gmail.com}\\[1.2em]
\today\\[1.6em]
%
\noindent
Clay Mathematics Institute\\
Scientific Advisory Board\\[1.8em]

\noindent Dear Members of the Board,\\[0.8em]

I enclose an article, \emph{Constructive Existence and Qualitative Uniqueness of Turbulent Navier--Stokes Solutions}, which I have
posted to arXiv. This letter says briefly what it contains.

Alternative~(A) of the Millennium Problem asks, for all admissible data and with no forcing, for a
solution of the incompressible Navier--Stokes equations that exists and is smooth for all time. What
the article presents is a constructive solution \emph{directed to}~(A): for data given in advance ---
a domain in three dimensions, initial data, forcing and boundary conditions, at high Reynolds number
--- it exhibits turbulent solutions and establishes that their mean-value outputs, such as drag and
lift, are determined to a stated precision. I do not claim this discharges~(A) as that alternative is
formulated, and the difference is stated plainly in the article: what is delivered is constructive
existence and qualitative uniqueness, and not smoothness, because the solutions that occur are not
smooth. Their gradients grow like $\nu^{-1/2}$ while their velocities stay bounded.

The method is the automated solution of differential equations, implemented in the FEniCS software:
a finite element discretisation of the given equation on a given mesh, with output error control from
the residual weighted by the solution of a dual problem, feeding back to optimisation of the mesh.
The error in each output is thus bounded by a quantity the computation itself produces.

Fefferman writes that the four alternatives are offered ``to give reasonable leeway to solvers while
retaining the heart of the problem''. The Institute's own description of the problem locates that
heart in turbulence. I have taken the leeway and used it to reach turbulence, and I offer the result
for what it is worth.

\paragraph{On prior publication.} The material is not new. The Clay problem was first addressed on
these lines in Chapter~14 of \emph{Computational Turbulent Incompressible Flow}, a book by Johan
Hoffman and myself published by Springer in 2007 as volume~4 of \emph{Applied Mathematics: Body and
Soul}, and the underlying computational record appears in refereed journals from 2006 onward, cited
in the article. The requirement of publication in a refereed journal, with two years elapsed, is thus
met by a margin approaching two decades. The remaining condition, that the work enjoy general
acceptance in the community, is not mine to declare.

\paragraph{On the occasion.} What is new is that a solution of alternative~(C) has been announced, by
methods themselves worth attention: a large optimisation fitting parameters to data by forward and
backward propagation, constructing an object rather than asserting that one exists. The article
observes that the method above is a restricted form of the same thing, and asks what it returns when
directed at~(A). It also records a scaling test showing that a singular solution of the kind
constructed for~(C) does not survive the limit of small viscosity in which fluid mechanics lies.

One interest should be declared: the industrial simulations cited in the article are the work of
Icarus Digital Math, of which I am a co-founder. The comparisons themselves are set and held by the
AIAA High Lift Prediction Workshop and the high-order CFD workshops, where predictions are submitted
before the measurements are disclosed.\\[1.4em]

\noindent Yours faithfully,\\[2.2em]

\noindent Claes Johnson\\
\noindent Professor Emeritus of Applied Mathematics\\
\noindent KTH Royal Institute of Technology

\end{document}
